\documentclass[11pt,a4paper]{article}
\usepackage[T1]{fontenc}
\usepackage[utf8]{inputenc}
\usepackage{lmodern}
\usepackage[margin=25mm,headheight=15pt]{geometry}
\usepackage{amsmath,amssymb,amsthm,mathtools}
\usepackage{mathrsfs}
\usepackage{booktabs,array,longtable}
\usepackage{microtype}
\usepackage{xcolor}
\usepackage{enumitem}
\usepackage{fancyhdr}
\usepackage{hyperref}
\hypersetup{colorlinks=true,linkcolor=blue!45!black,urlcolor=blue!55!black,
 citecolor=blue!45!black,pdftitle={Quadratic batching, index absorption, and smaller consistency equations},
 pdfauthor={Research synthesis prepared for ProveIt}}
\newcommand{\N}{\mathbb N}
\newcommand{\Z}{\mathbb Z}
\newcommand{\Q}{\mathbb Q}
\newcommand{\Con}{\operatorname{Con}}
\newcommand{\bitlength}{\operatorname{bitlength}}
\newcommand{\pc}{\operatorname{pc}}
\newcommand{\Zpos}{\Z_{>0}}
\newcommand{\repoPin}{\texttt{22e75dd6b87b1b199bb08c778d7cf075ae5aae08}}
\newtheorem{theorem}{Theorem}[section]
\newtheorem{lemma}[theorem]{Lemma}
\newtheorem{proposition}[theorem]{Proposition}
\newtheorem{corollary}[theorem]{Corollary}
\theoremstyle{definition}
\newtheorem{definition}[theorem]{Definition}
\theoremstyle{remark}
\newtheorem{remark}[theorem]{Remark}
\numberwithin{equation}{section}
\setlist{nosep,leftmargin=*}
\setlength{\parskip}{3pt}
\setlength{\parindent}{0pt}
\setlength{\emergencystretch}{2em}
\pagestyle{fancy}
\fancyhf{}
\fancyhead[L]{\small ProveIt: Diophantine consistency forms}
\fancyhead[R]{\small 23 October 2026}
\fancyfoot[C]{\thepage}
\begin{document}
\hypersetup{pageanchor=false}
\begin{titlepage}
\vspace*{18mm}
{\small\sffamily\color{blue!45!black} RESEARCH SYNTHESIS FOR PROVEIT}\par
\vspace{8mm}
{\Huge\bfseries Quadratic batching,\par index absorption,\par and smaller consistency\par equations}\par
\vspace{8mm}
{\Large Explicit Diophantine forms whose unsolvability\par is equivalent in ZFC to $\Con(\mathrm{ZFC})$}\par
\vspace{8mm}
{\large Prepared for Vladimir Reshetnikov\par 23 October 2026}
\vspace{12mm}
\begin{abstract}
We give three compatible improvements over the Diophantine consistency
forms at the pinned ProveIt repository revision. First, we test its
22 quadratic residuals simultaneously, using a 32-element Sidon set,
instead of first forming a quartic sum of squares and then using
exponential positional weights. The Jones interface yields a polynomial
in nine nonnegative integers of exact degree $166\,177\,600$ and an
868-operation evaluator; the strictly positive version uses 864
operations. A companion through the existing Sun interface has ten
unrestricted integer unknowns, exact degree $449\,607\,016$, and 773
operations. Second, a short-range Pell-index argument absorbs an index
quotient into the existing norm carrier, giving eleven positive unknowns
at degree 1154, or a two-equation system of degrees 2 and 1152; a
further positive translation gives degree 1140. Third,
complementing whole tile labels replaces bitwise clauses and reduces the
fixed native coefficient payload by about $33\%$.

All constants have finite recipes. The native tuple and source size
constants also have materialized binary payloads.
The archive contains explicit division-free circuits, a complete proof
of the new reductions, dependency records, and independent finite
certificates. The mathematical equivalences use the inherited native
compiler and literal K215 calibration. The ten-integer companion also
uses Sun's 2026 preprint. This report does not claim an end-to-end new
proof-assistant theorem, global minimality, or a proof of consistency.
\end{abstract}
\vfill
{\small Source revision: \repoPin\par
Repository: \url{https://github.com/VladimirReshetnikov/ProveIt}}
\end{titlepage}
\pagenumbering{roman}
\hypersetup{pageanchor=true}
\tableofcontents
\clearpage
\pagenumbering{arabic}
\section{Results, comparison, and logical meaning}
\label{sec:overview}

There is no single useful ordering of ``simplest'' Diophantine forms.
The number of unknowns, degree, number of equations, arithmetic circuit
size, and size of fixed coefficients are distinct resources. We therefore
give explicit forms on several parts of that tradeoff, with a common
fully specified coefficient tuple.

Throughout, $\N=\{0,1,2,\ldots\}$ and $\Zpos=\{1,2,\ldots\}$.
All additional letters introduced in formulas are polynomial
abbreviations, not existentially quantified unknowns. A circuit operation
means one binary addition, subtraction, or multiplication. Squaring and
multiplication by a fixed coefficient count; shared registers and fixed
leaves do not. Powers in displayed formulas always have fixed
nonnegative integer exponents and are expanded into multiplication gates
in the delivered circuits.

\begin{theorem}[Explicit consistency forms]
\label{thm:main}
With the four fixed native ports specified in
Section~\ref{sec:complemented-compiler}, the archive defines polynomials
$F_9$, $F_9^+$, $F_{10}$, and $F_{11}$ such that ordinary mathematics
formalizable in ZFC proves
\begin{align*}
 \Con(\mathrm{ZFC})
 &\ \longleftrightarrow\ \neg\exists\mathbf x\in\N^9\ F_9(\mathbf x)=0,\\
 &\ \longleftrightarrow\ \neg\exists\mathbf x\in\Zpos^9\ F_9^+(\mathbf x)=0,\\
 &\ \longleftrightarrow\ \neg\exists\mathbf x\in\Z^{10}\ F_{10}(\mathbf x)=0,\\
 &\ \longleftrightarrow\ \neg\exists\mathbf x\in\Zpos^{11}\ F_{11}(\mathbf x)=0.
\end{align*}
The last line also has an equivalent two-equation form in the same
eleven unknowns. The first two lines use the parameterized Jones ratio
theorems. The third additionally uses the stated Lucas and interpolation
results of Sun~\cite{sun2026}. Every line uses the inherited ordinary
compiler and literal K215 calibration at the pinned revision~\cite{repo}.
\end{theorem}

\begin{table}[htbp]
\centering
\small
\begin{tabular}{@{}llllr@{}}
\toprule
Form & Domain & Unknowns & Exact degree(s) & Operations\\
\midrule
Quadratic batching & $\N$ & 9 & $166\,177\,600$ & 868\\
Positive translation & $\Zpos$ & 9 & $166\,177\,600$ & 864\\
Sun companion & $\Z$ & 10 & $449\,607\,016$ & 773\\
Index absorption & $\Zpos$ & 11 & $1154$ & 123\\
Two-equation form & $\Zpos$ & 11 & $2,1152$ & 119\\
Shifted two-equation form & $\Zpos$ & 11 & $2,1152$ & 118\\
Translated product & $\Zpos$ & 11 & $1140$ & 129\\
Translated shifted system & $\Zpos$ & 11 & $2,1138$ & 124\\
\bottomrule
\end{tabular}
\caption{New forms in the delivered archive. The certificate circuit for
the nine-variable degree proof is intentionally a different evaluator.}
\label{tab:results}
\end{table}

\subsection{What improves over the pinned repository}

The nine-variable degree previously reached
\[
 92324\cdot5^{32}+296
 =2\,149\,585\,634\,469\,985\,961\,914\,062\,796.
\]
The new degree is smaller by a factor of approximately
$1.294\cdot10^{19}$. The 864-operation positive circuit also improves
the separate historical 1473-operation positive circuit by $41.34\%$;
the degree-optimized package~37 circuit uses 1500 operations. These are
separate baseline comparisons, not an invented baseline combining the
best features of different evaluators.

The ten-integer baseline in package~31 has degree
\[
 126496\cdot5^{31}+27064
 =589\,042\,901\,992\,797\,851\,562\,527\,064
\]
and 1442 operations. Quadratic batching reduces the degree by
approximately $1.310\cdot10^{18}$ and the circuit cost by $46.39\%$.
The domain remains unrestricted integers.

The eleven-positive construction removes one supplied unknown from the
direct twelve-positive Pell-norm family. It trades degree and arithmetic
for that reduction: the repository still has a twelve-positive degree
1022 construction with 124 operations and a separate 108-operation
degree 1081 construction. The present work does not improve the
historical 67-operation single equation or the 59-operation system.
It also leaves intact the usefulness of the 31-positive quadratic system
and its quartic sum of squares when low degree is the primary goal.

The native ports now use $339{,}476$ raw bytes rather than $506{,}664$,
a $32.998\%$ reduction. Their largest bit length falls from
$1{,}016{,}946$ to $681{,}374$. This is a reduction in native compiler
constants; it is not a claim about the size of a fully expanded
$166$-million-degree polynomial.

\subsection{The three ideas}

The first change is to preserve the system's quadratic structure.
Multiplying a quadratic coefficient kernel by the square of a digit code
can select a residual by one coefficient. A Sidon set gives unambiguous
quadratic monomials with maximum position 799. Disjoint translated
blocks select all 22 residuals simultaneously. A single binary carry
test enforces every target coefficient to be zero, so there is no need
to take a sum of squares before entering the few-variable interface.
This reduces both the degree and the work needed to evaluate the source
code. It also removes the old enormous Jones ceiling constant: the
nine-variable source can use simply $b=\epsilon+2$ and
$B=\kappa b^2$.

The second change is arithmetic rather than combinatorial. The
eleven-variable norm initially recovers only a weakened divisibility
$E\mid H(k-R-1)$. Cancelling $H$ here would be wrong. Instead we
recover the main Pell index first, place the remaining index in an
interval shorter than $E/\gcd(E,H)$, and deduce the exact equality
$2n=R+1$. This restores the full divisibility and reconstructs the
deleted positive quotient. The proof explicitly avoids a predecessor
bound that required the deleted quotient.
A further positive translation of the same carrier lowers degree 1154
to 1140 at the cost of six operations; it is supplied as a separate
optional form.

The third change uses a whole-label complemented encoding. An occupied
label $x$ is represented by $2^\ell-1-x$. A selected tile contributes
$x+1$, so one guarded binary field checks the entire label by requiring
its low $\ell$ bits to vanish. This replaces one guard field per bit
by one per label while retaining the occupied high bit needed by the
native compiler. We recompute the complete coefficient tuple and a
compatible packed-rank controller; changing only a few constants would
not suffice.

\subsection{What the consistency claim means}

Here $\Con(\mathrm{ZFC})$ means that no finite syntactic ZFC refutation
exists, using the effective proof presentation calibrated in the
repository. The fixed K215 machine satisfies the inherited theorem
\[
 \text{K215 halts on blank input}
 \quad\longleftrightarrow\quad
 \text{a finite ZFC refutation exists}.
\]
The named repository endpoint is
\texttt{KEndpoint215.blank\_halts\_iff\_zfc\_refutation}.
The compiler and arithmetic reductions connect a halt with an integer
zero, in both directions. Negation gives Theorem~\ref{thm:main}.

This is a theorem about equivalence. It does not prove either side of
the consistency statement. In particular it creates no conflict with
the second incompleteness theorem. All witnesses are ordinary finite
integers; none is a variable exponent, an unspecified enumeration index,
or an oracle for consistency.

The archive separates ordinary proofs, formal source endpoints, and
finite checks. We have not rebuilt the entire upstream proof-assistant
development or formalized the new batching and absorption arguments in
its kernel. The checks establish the finite facts claimed in their
receipts; they do not establish an unbounded equivalence by testing.
There is also no claim that nine is minimal for this particular
consistency predicate, or that the Sidon set or operation counts are
globally optimal.

\section{Direct quadratic batching and a nine-variable equation}
\label{sec:nine}

\subsection{The exact positive quadratic source}

We start with the literal 31-variable source of package~30, retained
by the later report packages~\cite{repo}. Its strictly positive unknowns,
in fixed order, are
\[
\begin{gathered}
 J,F,\alpha,z,f,h,i,o,j,s,w,\tau,\eta,\zeta,y,\gamma,\\
 N,E,P,c,a,X,D,C_2,L,F_2,V,g,I,B_a,\Delta.
\end{gathered}
\]
Only in this subsection, put $q=\beta J+1$, $k=\eta+\zeta$,
$H=4a+3$, and $C_w=q-F-\alpha$. The fixed ports are
$(\beta,K_0,M_C,M_F)$ from Section~\ref{sec:complemented-compiler}.
The 22 residuals, each required to vanish, are
\begin{alignat}{2}
 &L-q^2+1, &&N-q(L+1),\nonumber\\
 &X-wq, &&a-E-sN,\nonumber\\
 &E-X(a-E), &&c-k(a-E)-\eta,\nonumber\\
 &P-E(c-\eta)-k, &&\tau^2-P(P-k)-1,\nonumber\\
 &D-X-ac-\gamma H, &&\Delta-a^2-H,\nonumber\\
 &C_2-c^2, &&D^2-\Delta C_2-1,\nonumber\\
 &k-jc+V-1-hE, &&\label{eq:parent22}
\end{alignat}
\vspace{-6mm}
\begin{align*}
 &jc-V-(q-F)(N-q)+C_wL-(M_C+qM_F)J,\\
 &(K_0+w)C_w+q-F-z(q-1)-1,\\
 &F_2-f^2,\qquad V-of+c,\qquad g-y^2+V^2,\\
 &I-iC_2,\qquad B_a-I^2,\qquad B_a-\Delta(F_2-1),\\
 &y^2-B_ag-1.
\end{align*}
Let $r_0,\ldots,r_{21}\in\Z[n_1,\ldots,n_{31}]$ be these residuals
after replacing every positive source coordinate by $n_i+1$.
The delivered literal parent file \texttt{nine/parent31.json} has SHA-256
\begin{center}\footnotesize\ttfamily
1c336ba0b3c9827ffddcf7ce895cad9dca2907f907cde7484c7f387bb2cb1287.
\end{center}
Its historical binding metadata is preserved as provenance; the active
binding of this report is the new complete sibling compiler tuple.

For clarity, the elimination producing the 31-variable source is exact
on the positive domain. Its former register was
$G_p=(q-1)(q-F)+\alpha=q(q-F)-C_w$.
On $N=q(L+1)$ one has
\[
 LG_p=(q-F)(N-q)-C_wL.
\]
The transport row gives
$(K_0+w)C_w+(q-F)=1+z(q-1)\ge q$.
If $C_w\le0$, its left side is at most $q-F<q$, a contradiction.
Thus $C_w>0$, $q-F>0$, and the reconstructed $G_p$ is positive.
The unchanged normalized compiler theorem therefore gives
\begin{equation}
 \exists\mathbf n\in\N^{31}\ \bigwedge_{j=0}^{21}r_j(\mathbf n)=0
 \quad\longleftrightarrow\quad
 \text{K215 halts on blank input}.
 \label{eq:parent-calibration}
\end{equation}
The fifth residual at $\mathbf n=0$ is $1$, so a simultaneous zero
is never the all-zero vector. We use this small fact in completeness.

\subsection{A finite Sidon certificate replaces exponential positions}

Use the following increasing list of 32 positions:
\begin{equation}
\begin{split}
 (0,a_1,\ldots,a_{31})={}&(0,10,35,55,62,105,144,145,159,168,232,245,\\
 &306,323,325,353,371,374,382,494,499,525,566,578,\\
 &610,647,681,685,741,777,783,799).
\end{split}
\label{eq:weights}
\end{equation}
All 528 unordered pair sums with repetition are distinct. The file
\texttt{weights.json} and the independent verifier give a complete
finite certificate. No unverified extremality claim for this set is
needed. Put
\[
 M=799,\qquad d=4M+1=3197,\qquad \mathcal H=22d=70334.
\]
Write
\[
 r_j(\mathbf n)=\sum_{|\alpha|\le2}p_{j,\alpha}\mathbf n^\alpha,
 \qquad
 d_{j,\alpha}=(2-|\alpha|)!\prod_{i=1}^{31}\alpha_i!\,p_{j,\alpha}.
\]
Define the following fixed polynomials and fixed target positions:
\begin{align}
 E_j(T)&=\sum_{|\alpha|\le2}d_{j,\alpha}
                    T^{2M-\sum_i a_i\alpha_i},\nonumber\\
 E_*(T)&=\sum_{j=0}^{21}T^{jd}E_j(T),&
 Z_*(T)&=\sum_{j=0}^{21}T^{jd+2M}.
 \label{eq:kernel}
\end{align}
There are 136 nonzero weighted coefficient records and only 16
distinct nonliteral coefficient recipes. They are all explicit
polynomials in the four fixed native ports.

\begin{lemma}[Simultaneous quadratic extraction]
\label{lem:quadratic-extraction}
For $c(T)=1+\sum_{i=1}^{31}n_iT^{a_i}$,
\[
 [T^{jd+2M}]\,c(T)^2E_*(T)=2r_j(\mathbf n)
 \qquad(0\le j<22).
\]
\end{lemma}
\begin{proof}
Pad a monomial of total degree at most two with zero positions to
obtain a pair from the set in \eqref{eq:weights}. The pair-sum
certificate makes this representation unique. The coefficient of its
monomial in $c(T)^2$ is
$2!/((2-|\alpha|)!\prod_i\alpha_i!)$; this cancels the factorial
weight in $d_{j,\alpha}$, leaving $2p_{j,\alpha}\mathbf n^\alpha$.
Finally $c(T)^2E_j(T)$ has support in $[0,4M]$. Translating it by
$jd$ puts it in $[jd,jd+4M]$; these intervals are disjoint.
Thus no other row contributes to a target.
\end{proof}

The important distinction is between a polynomial's degree and the
choice of a positional code for its monomials. A quadratic system
needs unique pair sums, not separate base-five positions for a quartic
in 31 coordinates. Moreover the 22 target tests can share the same
unknown digit code; one does not need an extra unknown for each row.

\subsection{One explicit size constant}

Expand each weighted coefficient into monomials in the four positive
native ports. Let $C_*$ be one plus the sum of the absolute values of
all scalar coefficients times their corresponding port monomials,
over all 136 records. Thus every $|d_{j,\alpha}|<C_*$.
This is a bounded, finite integer expression, independent of the
existence of any witness. Set
\begin{equation}
 \kappa=2^{\bitlength(2048C_*)}>2048C_*.
 \label{eq:kappa}
\end{equation}
For the delivered tuple,
\[
 \bitlength(C_*)=1\,359\,514,\qquad
 \boxed{\kappa=2^{1\,359\,525}.}
\]
The constant's unsigned little-endian binary magnitude has 169,941
bytes; its compressed payload and all recipe hashes are included.
This replaces any unspecified coefficient-height parameter. The
derived fixed coefficients, including negative ones, remain signed
integer leaves of the final circuit.

\subsection{The new source code and its two-way interpretation}

For the moment let $b\ge2$ and $g>0$ be integers. Set
\begin{align}
 B&=\kappa b^2,&Q_0&=B^{M+1},&
 \ell&=\sum_{i=1}^{31}B^{a_i},\nonumber\\
 \lambda&=\sum_{t=0}^{\mathcal H-1}B^t,&
 S_0&=2(1+g)^2E_*(B)+B\lambda,\nonumber\\
 S&=g+Q_0S_0,&
 T&=Q_0-1-(b-1)\ell+Q_0(B-2)Z_*(B),\nonumber\\
 N&=B^{\mathcal H+M+2}=B^{71135},&
 R&=S(N^2-N)+(T+1)(N^2-1),\nonumber\\
 G_{\rm cap}&=32bB^M.&&
 \label{eq:nine-source}
\end{align}
Everything in this display is an integer polynomial in $b,g$ with
fixed integer coefficients. In particular the finite sum defining
$\lambda$ introduces no division in the polynomial.

\begin{lemma}[Unconditional nonvanishing and conditional size]
\label{lem:source-size}
For every $b\ge2,g>0$, the integer $R$ is nonzero. If additionally
$g<4bB^M$, then $0\le S,T<N$ and $R\ge N^2-1\ge N\ge\max\{b,8\}$.
\end{lemma}
\begin{proof}
Since $Q_0\mid N$, reduction of \eqref{eq:nine-source} modulo $Q_0$
gives
\[
 R\equiv (b-1)\ell\pmod{Q_0}.
\]
The positions in $\ell$ are distinct and at most $M$, while $b<B$.
Consequently $0<(b-1)\ell<Q_0$, proving $R\ne0$ without a sign
assumption on $E_*(B)$ or $R$.

Now suppose $g<4bB^M$. We have $g<Q_0$, because $B>4b$.
The highest exponent in $E_*$ is at most $\mathcal H-2M-1$.
Using $B\ge2$ and the coefficient bound,
\[
 |E_*(B)|<2C_*B^{\mathcal H-2M-1},\qquad
 1+g<5bB^M.
\]
Thus
\[
 |2(1+g)^2E_*(B)|
 <100C_*b^2B^{\mathcal H-1}<B^{\mathcal H}.
\]
Since $B^{\mathcal H}\le B\lambda<2B^{\mathcal H}$, it follows
that $0<S_0<3B^{\mathcal H}$. The low mask
$Q_0-1-(b-1)\ell$ is between $0$ and $Q_0-1$: its base-$B$ digits
are $B-b$ at the chosen positions and $B-1$ elsewhere. Also
$(B-2)Z_*(B)<B^{\mathcal H}$. The extra factor of $B$ in $N$
therefore gives $0\le S,T<N$. Finally $T+1\ge1$ and $S\ge0$ imply
$R\ge N^2-1$; the remaining bounds are immediate.
\end{proof}

Write $\tau_2(u,v)=0$ for addition without a binary carry, equivalently
disjoint sets of occupied binary positions for nonnegative $u,v$.

\begin{lemma}[All residuals in one carry test]
\label{lem:carry-decoder}
Suppose $b$ is a power of two, $b\ge2$, and $0<g<4bB^M$. Then
\[
 \tau_2(S,T)=0
 \quad\longleftrightarrow\quad
 \begin{cases}
 g=\sum_i n_iB^{a_i},\quad 0\le n_i<b,\\
 r_0(\mathbf n)=\cdots=r_{21}(\mathbf n)=0.
 \end{cases}
\]
Conversely, every simultaneous natural zero can be encoded in this
way for arbitrarily large powers of two $b$, with $0<g<bB^M$.
\end{lemma}
\begin{proof}
Now $B$ and $Q_0$ are powers of two. The strict block bounds in
Lemma~\ref{lem:source-size} split the carry condition into
\[
 \tau_2\bigl(g,Q_0-1-(b-1)\ell\bigr)=0,
 \qquad
 \tau_2\bigl(S_0,(B-2)Z_*(B)\bigr)=0.
\]
The first mask has all bits set at unchosen digit positions. At a
chosen position its digit $B-b$ has precisely the bits above the
lowest $\log_2b$ bits set. Hence it gives exactly the stated code
and digit bounds.

For decoded digits let $h_t$ be the coefficient of $T^t$ in
$c(T)^2E_*(T)$. Its degree is at most $\mathcal H-1$, and
\[
 |h_t|\le C_*\left(1+\sum_i n_i\right)^2
 <1024C_*b^2<B/2.
\]
Therefore the integer
\[
 S_0/2=(1+g)^2E_*(B)+(B/2)\lambda
\]
has ordinary base-$B$ digits $B/2+h_t$, all strictly between $0$
and $B$. Both summands in the second carry test are even. Divide
them by two; the mask becomes $(B/2-1)Z_*(B)$.
At each chosen target digit, $B/2-1$ occupies all the lower binary
positions. A digit strictly between $0$ and $B$ is disjoint from
this mask exactly when it equals $B/2$. By
Lemma~\ref{lem:quadratic-extraction}, the correction at the target
is $2r_j(\mathbf n)$, so every residual vanishes.

For completeness choose a power of two $b$ exceeding all coordinates
of a given zero, and use its digit code. Since the zero is not the
zero vector, $g>0$. The geometric estimate using $B>b$ gives
$g\le(b-1)\sum_{e=1}^M B^e<bB^M$. The same digit calculations,
now in the forward direction, prove both carry tests. The choice of
$b$ can be made arbitrarily large.
\end{proof}

The central-binomial carry lemma of Jones~\cite{jones1982}, provided
as \texttt{Jones1982.lemma\_2\_16} in the pinned
\texttt{Carries.lean}, states that for $N=2^a$ and $0\le S,T<N$,
\begin{equation}
 \tau_2(S,T)=0
 \quad\longleftrightarrow\quad
 N^2\mid\binom{2R}{R},
 \qquad R=S(N^2-N)+(T+1)(N^2-1).
 \label{eq:carry-binomial}
\end{equation}
This theorem uses precisely the source expression
\eqref{eq:nine-source}. It needs no quartic index or base-five
normalization. The source construction is now complete.

\subsection{The full nine-variable polynomial}

The nine nonnegative unknowns are
\[
 (\epsilon,g_0,h_0,s_0,w_0,\phi,i_0,j_0,n)\in\N^9.
\]
Put $b=\epsilon+2$ and
$g=g_0+1,h=h_0+1,s=s_0+1,w=w_0+1,i=i_0+1,j=j_0+1$.
Use $R,N,G_{\rm cap}$ from \eqref{eq:nine-source} and define
\begin{align}
 U&=N^2w,&Y&=N^2s,&P&=2(RY)^2U,\nonumber\\
 K&=R+1+h(P-1),&A&=RY(U+1),&B_1&=2R+1,\nonumber\\
 W_1&=bw,&C&=B_1+W_1+\phi,&J&=4A-5,\nonumber\\
 Q&=3W_1C-2(W_1^2-1),&D&=(A^2-1)C^2+1,\nonumber\\
 E&=iJDC^2,&F&=(A^2-1)E^2+1,\nonumber\\
 G_p&=A+F(CD+1-A),&H&=B_1+jC,&I&=(G_p^2-1)H^2+1,\nonumber\\
 a_0&=(P^2-1)K^2+1,&a_1&=DFI,&u&=JF,\nonumber\\
 v&=QF+J(H-C),&&&\nonumber\\[-1mm]
 m&=K^2G_{\rm cap}-4(C-KY)^2G_{\rm cap}-8gK^2.&&&
 \label{eq:jones-ratio}
\end{align}
Here $m$ is a register; the ninth unknown is $n$. Next set
\begin{align}
 \sigma&=Y(U+1),&r&=W_1+\phi,\nonumber\\
 M_D&=AC+\sigma(r+1)^2+1,&M_F&=AE+1,\nonumber\\
 T_{CH}&=CH+r^2,&
 M_I&=F(DT_{CH}+2H^2+A^2)-A^2+1,\nonumber\\
 V_0&=PK+1,&V_1&=M_DM_FM_I+1,\nonumber\\
 M_{\rm div}&=Q^2+(H-C)^2+1,&
 \Omega&=u(2m-1),\nonumber\\
 \Theta&=u(n+M_{\rm div})+v
          -\Omega(2M_{\rm div}+V_0V_1).&&
 \label{eq:majorants}
\end{align}
The required equation is
\begin{equation}
 \boxed{F_9=
   \bigl(\Theta^2-\Omega^2(a_0+V_0^2a_1)\bigr)^2
    -4\Omega^4V_0^2a_0a_1=0.}
 \label{eq:F9}
\end{equation}
Equations \eqref{eq:weights}--\eqref{eq:F9}, together with the four
native ports, specify a literal integer-coefficient polynomial. The
JSON arithmetic circuit is its expanded straight-line definition.
No square roots, tests, divisions, or variable exponents occur in it.

\subsection{Signed bounds before positive source recovery}

It is essential not to assume $R>0$ when applying the final polynomial.
For a freely chosen large $g$, $E_*(B)$ can make $R$ negative. The
following proof makes the order of implications explicit.

By Lemma~\ref{lem:source-size}, $R\ne0$ for every assignment of
the eight outer natural coordinates. Since $U,Y\ge1$,
\[
 |A|\ge2,\quad P\ge2R^2\ge2,\quad
 K\ge R+P\ge2R^2+R\ge1.
\]
Consequently $J\ne0$, $AJ>0$, $D\ge1$, $AE\ge0$, and $F\ge1$.
The completed-square bounds from package~37 apply with no sign
restriction on $R$. For completeness, here are their short proofs.
Because $A=R\sigma$ and $C=2R+r+1$,
\begin{align*}
 M_D-AC&=\sigma(r+1)^2+1>0,\\
 M_D+AC&=\sigma\bigl(4R^2+2R(r+1)+(r+1)^2\bigr)+1>0.
\end{align*}
The quadratic expression is nonnegative; twice it is
$(2R+r+1)^2+4R^2+(r+1)^2$.
Thus $M_D\ge|AC|+1$ and
$D=(AC)^2-C^2+1\le M_D^2$.
Similarly $M_F=AE+1\ge1$ and $F\le M_F^2$.

Since $H=(j+1)C-r$,
\[
 T_{CH}-CH=r^2\ge0,\qquad
 T_{CH}+CH=C^2+(C-r)^2+2jC^2\ge0.
\]
It follows that $T_{CH}\ge|CH|$.
Using $G_p=FDC+F-A(F-1)$, the integer inequalities
$|H|\le H^2$ and $|AH|\le A^2+H^2$ imply
\[
 |G_pH|+|H|+1
 \le FDT_{CH}+2FH^2+(F-1)A^2+1=M_I.
\]
Thus $M_I\ge1$ and $\sqrt{|I|}\le M_I$, even when the candidate
$I$ is negative. Altogether
\begin{equation}
 V_0,V_1\ge2,\qquad
 \sqrt{|a_t|}\le V_t-1\quad(t=0,1).
 \label{eq:root-bounds}
\end{equation}
Indeed $1\le a_0\le(PK)^2$ and
$\sqrt{|a_1|}\le M_DM_FM_I$. Also $u\ne0$ and
\begin{equation}
 \left|\frac vu\right|
 =\left|\frac QJ+\frac{H-C}{F}\right|
 \le|Q|+|H-C|\le M_{\rm div}-1.
 \label{eq:quotient-bound}
\end{equation}
All these statements precede square conditions, divisibility, and $m>0$.

\begin{lemma}[Signed-divisor finalizer]
\label{lem:finalizer}
Let integers $a_0,a_1,u,v,m,M,V_0,V_1$ satisfy $u\ne0$, $M\ge1$,
$|v/u|\le M-1$, $V_0,V_1\ge2$, and
$\sqrt{|a_t|}\le V_t-1$. Define $\Omega$, $\Theta$, and $F_9$ as
above, with $M_{\rm div}=M$. Then
\[
 \exists n\in\N\ F_9=0
 \quad\longleftrightarrow\quad
 a_0,a_1\text{ are integer squares},\quad u\mid v,\quad m>0.
\]
\end{lemma}
\begin{proof}
The odd integer $2m-1$ is nonzero, so $\Omega\ne0$. Choose complex
roots $r_t^2=a_t$. Direct multiplication factors $F_9$ into
\[
 \prod_{\varepsilon,\delta\in\{-1,1\}}
       \bigl(\Theta-\Omega(\varepsilon r_0+\delta V_0r_1)\bigr).
\]
A zero gives one vanishing factor. Absorb its signs in the chosen
roots, and put $s=\Theta/\Omega=r_0+V_0r_1\in\Q$.
If $s\ne0$, the identity
$2sr_0=s^2+a_0-V_0^2a_1$ makes both roots rational. A rational
number with integer square is an integer. If $s=0$, the bounds give
$|r_1|=|r_0|/V_0<1$, hence the integer $|a_1|<1$; then
$a_1=a_0=0$. This includes negative candidate radicands and coincident
square classes, without presuming independent field automorphisms.

Set $Z=2M+V_0V_1+r_0+V_0r_1$. The root bounds give
$Z\ge2M+1>0$. Dividing the vanishing factor by the possibly negative
$u$ gives
\[
 n+M+v/u=(2m-1)Z.
\]
Its left side is at least $n+1>0$, so $m>0$. Its right side is
integral, so $u\mid v$. Conversely, if these three conditions hold,
choose integer roots and define
$n=(2m-1)Z-M-v/u$. This is an integer at least $2$, and makes
the corresponding factor vanish.
\end{proof}

\subsection{Soundness and completeness of the nine-variable form}

The parameterized ratio theorems in
\texttt{NineRatio.lean} have the following precise interface.
For $R,N\ge8$, $0<b\le\min\{R,N\}$, and $G>0$, the conditions
\[
 a_0,a_1\text{ squares},\qquad JF\mid QF+J(H-C),\qquad m>0
\]
with $h,s,w,i,j>0$ and $\phi\ge0$ imply
\begin{equation}
 N^2\mid\binom{2R}{R},\qquad b\text{ is a power of two},
 \qquad8g<G.
 \label{eq:ratio-sound}
\end{equation}
Conversely, the first two conditions in \eqref{eq:ratio-sound}
together with $32g<G$ produce positive $h,s,w,i,j$ and
nonnegative $\phi$ satisfying the three displayed ratio conditions.
These are \texttt{Jones1982.Nine.ratio\_sound} and
\texttt{ratio\_complete}; the formulas match \eqref{eq:jones-ratio}
with $G=G_{\rm cap}$. They quantify arbitrary admissible $R,N,b,G$
and do not require Jones's original source packing.

Suppose $F_9=0$. The unconditional bounds and
Lemma~\ref{lem:finalizer} first give $m>0$. As $K\ge1$ and the
subtracted square term in $m$ is nonnegative, this already implies
\[
 8g<G_{\rm cap},\qquad g<4bB^M.
\]
Only now apply Lemma~\ref{lem:source-size}: it supplies positive
$R,N$ and every size hypothesis of the ratio theorem. Apply
\eqref{eq:ratio-sound}, then \eqref{eq:carry-binomial}, then
Lemma~\ref{lem:carry-decoder}. This gives a zero of the 22 parent
residuals, and hence a blank K215 halt by
\eqref{eq:parent-calibration}.

Conversely suppose the parent residuals have a natural zero. Choose
$b$ and its positive code $g<bB^M$ by
Lemma~\ref{lem:carry-decoder}. Its carry condition and
\eqref{eq:carry-binomial} give the binomial divisibility; all ratio
size conditions hold. Since $32g<32bB^M=G_{\rm cap}$,
\texttt{ratio\_complete} supplies the six remaining outer witnesses.
Lemma~\ref{lem:finalizer} supplies $n\in\N$. Translating the positive
outer witnesses down by one gives a zero of $F_9$. We have proved
\[
 \exists\N^9\ F_9=0
 \quad\longleftrightarrow\quad
 \exists\N^{31}\ \bigwedge_jr_j=0,
\]
and thus the first line of Theorem~\ref{thm:main}.

For strictly positive coordinates use
\[
 (\epsilon_+,g,h,s,w,\phi_+,i,j,n_+)\in\Zpos^9,
 \quad b=\epsilon_++1,\quad\phi=\phi_+-1,\quad n=n_+-1.
\]
This coordinate bijection gives $F_9^+$, preserves exact degree,
and saves four operations in the delivered evaluator. It does not
extend the result to nine unrestricted integers.

\subsection{Exact degree, with cancellations resolved}

Treat fixed coefficients as degree zero. The actual degree cannot
be read from the original nested expression by blindly taking
maxima: its highest apparent terms cancel. The following exact
identities expose the surviving leading term. Put
$t_0=\sigma(r+1)^2$ and
\begin{align*}
 \delta_D&=M_D^2-D=2ACt_0+t_0^2+2AC+2t_0+C^2,\\
 \delta_F&=M_F^2-F=2AE+E^2,\\
 W_-&=M_I-G_pH\\
 &=FDr^2+2FH^2+(F-1)A^2+1-FH+A(F-1)H,\\
 \delta_I&=M_I^2-I=W_-(M_I+G_pH)+H^2-1,\\
 \delta_V&=V_1^2-DFI\\
 &=\delta_DM_F^2M_I^2+D\delta_FM_I^2+DF\delta_I
                      +2M_DM_FM_I+1.
\end{align*}
With $E_0=u(n+M_{\rm div})+v-2\Omega M_{\rm div}$, the inner
square in \eqref{eq:F9} is exactly
\begin{equation}
 E_0^2-2E_0\Omega V_0V_1+\Omega^2V_0^2\delta_V-\Omega^2a_0.
 \label{eq:degree-inner}
\end{equation}
Some decisive total degrees are as follows:
\[
\begin{array}{c|r@{\qquad}c|r}
 N&142270&R&426808\\
 P&1707239&K&1707240\\
 A&995890&C&426808\\
 D&2845396&F&11381586\\
 I&30161198&a_1&44388180\\
 V_0&3414479&V_1&22194090\\
 \Omega&16362637&\delta_D&1991784\\
 \delta_F&9389806&\delta_I&29307585\\
 \delta_V&43534568&E_0&17216257
\end{array}
\]
The first three terms defining $\delta_V$ have degrees
$43\,534\,568$, $42\,396\,400$, and $43\,534\,567$.
The first uniquely dominates. The third term in
\eqref{eq:degree-inner} uniquely dominates, at degree $83\,088\,800$.
The final subtracted term of \eqref{eq:F9} has strictly smaller degree
than twice this value. Hence
\begin{equation}
 \boxed{\deg F_9=166\,177\,600.}
 \label{eq:degree9}
\end{equation}
For a nonvanishing certificate, substitute all nine natural inputs by
$t$. The coefficient of $t^{166177600}$ is exactly
\[
 \boxed{2^{196}\kappa^{83088652}>0.}
\]
The verifier tracks at every gate a total-degree upper bound and the
exact coefficient of the corresponding diagonal power. All sixteen
derived native coefficients remain independent symbolic leaves; none
is assigned zero or presumed generic. The final coefficient is
independent of them. This proves the claimed degree for the actual
materialized tuple as well as for every compatible specialization.

The production circuits use sparse Horner evaluation of $E_*$,
shared repeated squaring, and a divide-and-conquer geometric-sum
identity. Their literal gate counts are 868 for $\N^9$ and 864 for
$\Zpos^9$. The separate cancellation certificate uses 917 operations;
it is a proof artifact, not the claimed fastest evaluator.

\section{A companion equation in ten unrestricted integers}
\label{sec:sun-companion}

The same quadratic coefficient packing also compresses the ten-integer
construction already present in source~31 of the repository~\cite{repo}.
The resulting polynomial has
\[
 \boxed{10\text{ integer unknowns},\qquad
 \deg F_{10}=449\,607\,016,\qquad 773\text{ binary operations}.}
\]
The earlier source~31 equation has degree
$589\,042\,901\,992\,797\,851\,562\,527\,064$ and 1,442 operations.
The comparison uses the same convention: every binary addition,
subtraction and multiplication is charged, including multiplication by
a fixed coefficient; repeated registers and fixed leaves are free.
Thus this application improves both degree and circuit size, while
retaining the same ten unrestricted integer unknowns.  It adds the
Lucas-theoretic dependency on Sun's 2026 preprint~\cite{sun2026}; the
nine-natural-variable construction does not need that preprint.

\subsection{A different source scale, with the same quadratic kernel}

Use the same 31 shifted quadratic coordinates and 22 residuals $r_j$,
the same Sidon positions $0,a_1,\ldots,a_{31}$, and the same
$M=799$, $d=3197$, $H=70334$ as in the preceding packing construction.
With $|\alpha|\le2$, write
\[
 r_j(\mathbf n)=\sum_\alpha p_{j,\alpha}\mathbf n^\alpha,
 \qquad c_{j,\alpha}=(2-|\alpha|)!\prod_i\alpha_i!\,p_{j,\alpha}.
\]
Define the fixed polynomial and target mask
\begin{align*}
 E(T)&=\sum_{j=0}^{21}\sum_{|\alpha|\le2}
 c_{j,\alpha}T^{jd+2M-\sum_i\alpha_i a_i},\\
 Z(T)&=\sum_{j=0}^{21}T^{jd+2M}.
\end{align*}
There are 136 nonzero weighted coefficient records and 16 distinct
nonliteral coefficient recipes. All four native ports are bound to the
new tuple compiled from complemented labels in the sibling \texttt{compiler} directory.

For each monomial $q\prod_i\mathrm{port}_i^{e_i}$ in these records,
majorize its magnitude by $|q|2^{\sum_i e_i b_i}$, where $b_i$ is the
actual bit length of the corresponding native port. Let $C_{\mathrm S}$ be one
plus the sum of all these majorants. This bit-length majorant is distinct
from the exact positive-monomial majorant used for the nine-variable source. The delivered exact arithmetic gives
\[
 \operatorname{bitlength}(C_{\mathrm S})=1\,359\,515.
\]
Set
\[
 \beta_{\mathrm S}=2^{1\,359\,528}>2048C_{\mathrm S}.
\]
Its exponent is even, so $\beta_{\mathrm S}\equiv1\pmod3$.
The binary payload of this additional constant occupies 169,942 bytes
before compression and 200 bytes in the delivered compressed form.
The number is fixed by the finite compiler data, independently of
whether a refutation exists.

For integer variables $f,g$, define
\begin{align*}
 b&=1+3f,& \mathcal B&=\beta_{\mathrm S}b^2,&
 \ell&=\sum_{i=1}^{31}\mathcal B^{a_i},\\
 \lambda&=\sum_{t=0}^{H-1}\mathcal B^t,&
 N_0&=\mathcal B^{M+1},& N_1&=4\mathcal B^H,\\
 N&=N_0N_1=4\mathcal B^{71134},&
 \mathscr C&=b\mathcal B^M,&
 M_0&=N_0-1-(b-1)\ell.
\end{align*}
Put
\[
 J_0=2(1+g)^2E(\mathcal B)+\mathcal B\lambda,
\]
and then
\begin{align*}
 S_0&=g+N_0J_0,&
 T_0&=M_0+N_0(\mathcal B-2)Z(\mathcal B),\\
 R_0&=(S_0+T_0+1)N+T_0+1,&
 X&=(N-1)R_0,&Y&=N^2.
\end{align*}
Here $E$ is a polynomial abbreviation rather than a new unknown.
All powers and summation limits are fixed. No sign restriction is
imposed on these polynomials for arbitrary free inputs: the reverse
bridge first recovers the bounds used in the next paragraph.

\paragraph{The source contract.}
These formulas have the following three properties.
If the quadratic system has a natural solution, there are arbitrarily
large $f>0$ and a $g$ with $b\le g<\mathscr C$, such that $b$ is a
square and a power of two, and $Y\mid\binom{2X}{X}$.
Conversely, $f\ne0$, $0\le g<2\mathscr C$, the same two conditions on
$b$, and that divisibility imply a natural solution of the system.
Finally, whenever $f\ne0$, $b$ is a square and $0\le g<2\mathscr C$,
one has $X>3b$ and $Y>\max\{b,256\}$, even before imposing the
power-of-two condition.

\paragraph{Proof of the changed source step.}
A square $b=1+3f$ with $f\ne0$ has $b\ge4$ and $f>0$.
The highest exponent of $E$ is at most $H-2M-1$, so
\[
 |E(\mathcal B)|<2C_{\mathrm S}\mathcal B^{H-2M-1}.
\]
For $0\le g<2b\mathcal B^M$, the inequality
$1+g\le3b\mathcal B^M$ gives
\[
 |2(1+g)^2E(\mathcal B)|
 <36C_{\mathrm S}b^2\mathcal B^{H-1}<\mathcal B^H.
\]
Since $\mathcal B^H\le\mathcal B\lambda<2\mathcal B^H$,
we obtain $0<J_0<3\mathcal B^H<N_1$ and $g<N_0$.
The mask $M_0$ has digits $\mathcal B-b$ at the positions $a_i$
and $\mathcal B-1$ elsewhere, so $0\le M_0<N_0$.
Also $(\mathcal B-2)Z(\mathcal B)<N_1$.
Consequently $0\le S_0,T_0<N$, $R_0\ge N+1$,
$X\ge N^2-1>3b$, and $Y>\max\{b,256\}$.
Moreover $N\equiv1\pmod3$, so $3\mid X$.

If $b$ is a power of two, so are $\mathcal B,N_0,N_1,N$.
Sun's central-binomial carry lemma~\cite[Lemma 2.2]{sun2021} gives
\[
 Y\mid\binom{2X}{X}
 \quad\Longleftrightarrow\quad
 \tau_2(S_0,T_0)=0.
\]
The two blocks split this into the code-validity test
$\tau_2(g,M_0)=0$ and the residual test
$\tau_2(J_0,(\mathcal B-2)Z(\mathcal B))=0$.
The first test decodes
$g=\sum_i n_i\mathcal B^{a_i}$ with $0\le n_i<b$ by
\cite[Lemma 2.4]{sun2021}, which explicitly permits arbitrary distinct
digit positions. For such digits, put
\[
 K_0=(1+g)^2E(\mathcal B)+(\mathcal B/2)\lambda,
 \qquad J_0=2K_0.
\]
Every coefficient of $K_0$ lies strictly between 0 and $\mathcal B$,
because its signed correction has magnitude at most
\[
 C_{\mathrm S}\left(1+\sum_i n_i\right)^2
 <1024C_{\mathrm S}b^2<\mathcal B/2.
\]
The selected digits are $\mathcal B/2+2r_j(\mathbf n)$.
Adding $(\mathcal B/2-1)Z(\mathcal B)$ without a binary carry
forces each selected digit to equal $\mathcal B/2$, hence every
$r_j(\mathbf n)=0$. This proves soundness.
Conversely, a simultaneous zero is nonzero because the fifth residual
at $\mathbf n=0$ equals 1. Choose $b=4^q$ larger than all its coordinates,
set $f=(b-1)/3$, and encode it in $g$. Then $b\le g<\mathscr C$,
and both carry tests hold. This proves completeness, with $q$ arbitrarily
large. No sum of squares is used in this source step.

\subsection{The precise imported Lucas interface}

The positive and reverse bridge proofs in
\cite[Propositions 6.1 and 6.2]{sun2026} use precisely the three-part
source contract of its Proposition 2.1, established above for the new
source, the fixed form $b=1+3f$, and the monomial bound
$\mathscr C=2^{(1\,359\,528)M}b^{2M+1}$.
Their proofs do not otherwise use the old base-five coefficient code.
We use this proof-level specialization, rather than asserting that the
new source is literally Sun's original source polynomial.
The positive bridge supplies a Lucas certificate, one square and one
divisibility condition; the reverse bridge recovers the source contract
from those conditions and a positivity inequality.
Sun's modified Lucas certificate~\cite[Theorem 3.1]{sun2026}
and full residue interpolation~\cite[Theorem 4.1]{sun2026} are additional
imported results. The latter permits the second Lucas witness to lie
in any prescribed residue class modulo a suitable first witness.
This permits a positivity witness and the divisibility
merger to share variables. The positive bridge is the specialization
$p=P=2$, $Q=1$ of \cite[Theorem 4.12]{sun2021}; the reverse bridge
uses the estimates of \cite[Theorem 4.13]{sun2021}, replacing its
Lucas-certificate call by \cite[Theorem 3.1]{sun2026}.

The complete ten-variable list is
\[
 (f,g,h,k,l,w,x,y,z,m)\in\Z^{10}.
\]
Use the source quantities above and define
\begin{align*}
 L&=lY,&U&=2LX,\\
 V&=4gwY,&W&=bw,\\
 A&=U(V+1),&B_1&=2X+1,\\
 C_1&=B_1+(A-2)h,&P&=U^2V,\\
 K&=X+1+k(P-2),&S&=2A-5,\\
 T&=3WC_1-2(W^2-1),&r_1&=(P^2-4)K^2+4.
\end{align*}
Next set
\begin{align*}
 \Delta&=8\mathscr C^3gK^2
       -g^2\bigl(32(C_1-KL)^2\mathscr C^3+g^2K^2\bigr),\\
 O&=f^2l^2\Delta,\qquad\xi=zx.
\end{align*}
The modified Lucas quantities are
\begin{align*}
 D&=(A^2-4)C_1^2+4,& E_1&=C_1^2D\xi,\\
 F&=4(A^2-4)E_1^2+1,&G&=2(A+2)(A-2)^2E_1^2-1,\\
 \Omega_1&=16E_1,&H_1&=C_1+(B_1-C_1)F+\Omega_1Fy,\\
 I&=(G^2-1)H_1^2+1,&r_2&=DFI,\\
 r_0&=b\bigl(1+(3O-4)b^2z^2\bigr).&&
\end{align*}
The positive bridge has $O\ge4$, since $g\ge b\ge4$ and its positive
integer $\Delta$ is divisible by $g$.
As proved in source~31~\cite{repo}, the coprimality of $b$ and
$1+(3O-4)b^2z^2$, together with $b\equiv1\pmod3$, makes $r_0$ square
and $z\ne0$ enforce $b$ square and $O>0$. Conversely, a Pell solution
with its second coordinate divisible by $b$ supplies such $z$.

\subsection{The final polynomial and its implication in both directions}

Let
\begin{align*}
 \pi(v)&=(v^3-v-1)(v^2+23),&\Sigma&=\xi S^5,\\
 t&=T+Sy,&
 \Theta&=(t^3-tS^2-S^3)(t^2+23S^2).
\end{align*}
Define the weights and scaled radicands by
\[
 w_0=1,\qquad w_1=r_0,\qquad w_2=r_0(P^2+K^2+3),
\]
\[
 Z_0=\Sigma m-\Theta,\qquad
 A_i=\Sigma^2w_i^2r_i\quad(0\le i\le2),
\]
and put
\[
 V_0=Z_0^2+A_2-A_0-A_1,\qquad
 W_0=V_0^2+4Z_0^2A_2-4A_0A_1.
\]
The requested integer equation is
\[
 \boxed{F_{10}=W_0^2-16V_0^2Z_0^2A_2=0.}
\]
Every symbol other than the ten listed inputs is a polynomial
abbreviation or a fixed integer.

The local quintic merger and norm arguments are inherited from
source~31~\cite{repo}. Byte-exact copies of that proof, its mathematical
review and its finalizer are supplied under \texttt{signed10/inherited/};
the historical coefficient binding there is inactive. For clarity, their exact roles are as follows.
The polynomial $\pi$ has no rational root and has a root modulo every
positive integer. Thus $\Theta\ne0$, since $S=2A-5\ne0$.
Divisibility $\Sigma\mid\Theta$ forces $\xi\ne0$ and $S\mid T$:
if $v_p(S)>v_p(T)$, the two homogeneous factors have valuations
$3v_p(T)$ and $2v_p(T)$, contradicting divisibility by $S^5$.
Conversely, if $S\mid T$ and $\xi\ne0$, a suitable residue class of $y$
modulo $|\xi|$ makes the divisibility true. Full residue interpolation
selects that class while preserving the Lucas square certificate.
This converse is existential in $y$, not a pointwise assertion for every $y$.

The norm identity is
\[
 F_{10}=\prod_{\varepsilon\in\{\pm1\}^3}
 \left(Z_0-\Sigma\sum_{i=0}^2\varepsilon_iw_i\sqrt{r_i}\right).
\]
At a zero, $\Sigma=0$ would imply $F_{10}=\Theta^8\ne0$, so
$z\ne0$. Then $|r_0|>1$, and the structural weights satisfy
\[
 |w_1|>\sqrt{|r_0|},\qquad
 |w_2|>\sqrt{|r_0|}+|w_1|\sqrt{|r_1|}.
\]
A rational signed sum of these roots must have each root rational:
a Galois automorphism changing the largest indexed root cannot cancel
its change with the smaller terms. Rational square roots of integers
are integers, so the three radicands are squares and $\Sigma\mid\Theta$.
This reasoning also covers negative candidate radicands, zero roots,
and repeated square classes.

The reverse Lucas bridge therefore gives a simultaneous natural zero
of the 22 residuals. Conversely, its positive bridge, the coprime Pell
step and full residue interpolation give the three square conditions
and $\Sigma\mid\Theta$; an integer $m$ then kills one norm factor.
Together with the native compiler and calibrated K215 endpoint, this
proves
\[
 \mathrm{ZFC}\vdash
 \bigl(\neg\exists\mathbf x\in\Z^{10}\ F_{10}(\mathbf x)=0\bigr)
 \longleftrightarrow\operatorname{Con}(\mathrm{ZFC}).
\]
This is an equivalence based on the stated mathematical dependencies,
not a proof of consistency or a new proof-assistant certificate.

\subsection{Degree and actual-instance checks}

The degrees of $X,Y$ are 426,804 and 284,536. The three scaled radicands
$A_0,A_1,A_2$ have degrees 13,947,118, 24,764,324 and 112,401,754.
The last dominates the squared norm center and the other radicands, so
the final leading part is $A_2^4$. The claimed exact degree is therefore
$4\cdot112\,401\,754=449\,607\,016$.
Independently, diagonal leading coefficients for the actual complemented
tuple are nonzero modulo 101, 1009 and 1,000,003: they are 81, 287 and 270,876.
These are exact finite certificates that the structural degree bound is
attained after specializing all fixed coefficients.

The verifier independently reconstructs the 136 records by evaluating
quadratic values, gradients and homogeneous parts at the all-one
coordinate vector, while the builder uses binomial expansion. It checks
all 528 pair sums, every native-port hash and strict coefficient bound,
the materialized $\beta_{\mathrm S}$, the formal norm expansion, and
1,620 named-register comparisons against literal scalar formulas over
three finite fields. These checks establish the finite artifact's
transcription and quantitative claims; the unbounded equivalence rests
on the proofs above and the stated imported theorems.

\section{Eleven positive unknowns: absorb the index quotient}
\label{sec:eleven}

The second new mechanism removes an actual supplied variable from the direct
Pell construction. It gives a single equation in eleven positive integers,
of exact degree $1154$ and cost $123$ binary operations, or two equations of
maximum degree $1152$ and cost $118$ after the inherited transport shift.
The predecessor direct family uses twelve positive integers and three
equations. This is a different tradeoff from the nine-variable Jones
construction: its advantage is a moderate degree and a small explicit circuit.
The comparison is with the pinned repository collection~\cite{repo}, not a
global minimality assertion.

The old index quotient $h$ is reconstructed from a divisibility condition
inside the already quantified carrier. The first divisibility recovered is
weaker than the original one. A short-range argument for the Pell index makes
it sufficient. This is why simply deleting the old index equation would not
have worked.

\subsection{The polynomial and its exact supplied domain}

Use the entire fixed coefficient tuple $(\beta,K_0,M_C,M_F)$ of the compiler,
with $M_F$ the shifted high mask. The eleven supplied unknowns are exactly
\[
 J,F,\alpha,z,i,T,s,w,\eta,\zeta,t\in\Z_{>0}.
\]
All other letters below are polynomial abbreviations:
\begin{align*}
 q&=\beta J+1,&s_0&=s+3q,&C_w&=q-F-\alpha,\\
 X&=wq,&Y&=s_0q^3,&E&=XY,\\
 k&=\eta+\zeta,&c&=kY+\eta,&a&=E+Y,\\
 H&=4a+3,&\Delta&=a^2+H,&\ell&=EkY,\\
 A&=1+\ell(\ell+k),&G&=1+\Delta c^2,&b_0&=X+ac,\\
 I&=ic^2,&D&=1+\Delta I^2,&Q&=\Delta(D-1)=(\Delta I)^2,\\
 d&=Q-1,&v&=cT,&b&=c+RD.
\end{align*}
Here $R$ and the transport expression are
\begin{align*}
 R&=((q-1)(q-F)+\alpha)(q^2-1)+(M_C+qM_F)J,\\
 N_t&=(K_0+w)C_w+q-F-z(q-1).
\end{align*}
Define the sharp weight and new carrier by
\begin{align*}
 e&=Q(Dv^2+b^2)-1,&B_a&=2Qvb,\\
 C_*&=E+2(X+Y)+s_0q^2,&M&=I(2T+RC_*ic),\\
 L_0&=1+HM,&S&=HE,\\
 L&=EL_0,&W&=Lb_0+H(Et+k-R-1).
\end{align*}
The finite conjugate product
\begin{equation}
 \mathcal N=
 \prod_{\substack{g=\pm\sqrt G,\ \tau=\pm\sqrt A\\f=\pm\sqrt D}}
 \left[d(W-Lg-SM\tau)^2-S^2e+S^2B_af\right]
 \label{eq:eleven-norm}
\end{equation}
is an integer polynomial. Its eight formal factors define a polynomial
identity; the displayed radicals are not additional unknowns, and independent
formal signs are not assumed to be independent field automorphisms.
The single equation and two-equation form are respectively
\begin{equation}
 \boxed{N_t(1+\mathcal N)-1=0},
 \qquad
 \boxed{N_t=1,\quad\mathcal N=0}.
 \label{eq:eleven-final}
\end{equation}

For a literal division-free evaluation, put
\begin{align*}
 n&=(SM)^2A,&d_n&=dn,\\
 P_H&=S^2QDv^2,&U_H&=S^2Qb^2,\\
 e_H&=P_H+U_H-S^2,&j&=4P_HU_H,\\
 \mathcal T&=2d(W^2+L^2G),&\kappa&=[d(W^2-L^2G)]^2,\\
 \mathcal A&=2(d_n+e_H),&\mathcal B&=(d_n-e_H)^2-j,\\
 a_r&=\mathcal T-\mathcal A,&b_r&=\mathcal B-\kappa,\\
 v_r&=16d_nj-a_r(a_r\mathcal T+2b_r),&u_r&=b_r^2-a_r^2\kappa,\\
 \mathcal N&=\kappa v_r^2-u_r(\mathcal T v_r-u_r).
\end{align*}
The system compares the two final terms directly. The identity follows by
reducing a quartic modulo the quadratic with root sum $\mathcal T$ and root
product $\kappa$, exactly as in the predecessor norm; it remains valid when
quadratic fields coincide or a radicand specializes to a square.

\subsection{Recover roots without assuming the deleted equation}

At a root of the product equation, $N_t$ and $1+\mathcal N$ are integer
units. For $x_0=d(W^2+L^2G)$, one has
\[
 \mathcal T=2x_0,\qquad
 \kappa=x_0^2-4d^2W^2L^2G,\qquad
 \mathcal N\equiv(u_r-x_0v_r)^2\pmod4.
\]
Thus $1+\mathcal N\ne-1$, and $\mathcal N=0$, $N_t=1$. The authentic
transport interface now gives
\begin{equation}
 q\ge16,\quad C_w>0,\quad 3q<R<q^4-2,\quad
 Y>3q^4>3R,\quad X\ge16,\quad E>R+1.
 \label{eq:eleven-outer}
\end{equation}
Also $k\ge2$ and $kY<c<k(Y+1)$. For example, $C_w\le0$ would imply
$N_t\le1-F-(z-1)(q-1)\le0$. We do \emph{not} assume $k>E$ or $R<k$:
those predecessor bounds depended on the omitted equation.

\begin{lemma}[Uniform auxiliary bound]
Every coherent complex auxiliary branch satisfying
\[
 f^2=D,\qquad dy^2=e-B_af=Q(vf-b)^2-1
\]
obeys $|y|<M\sqrt A$ under \eqref{eq:eleven-outer}.
\end{lemma}
\begin{proof}
Since $\Delta>R+2$ and $I\ge c^2$,
\[
 Q-b=\Delta I^2(\Delta-R)-c-R>0,
 \qquad \sqrt D<\Delta I<d.
\]
For $Z=v\sqrt D+b>1$,
\[
 |y|^2\le\frac{QZ^2+1}{d}
 <[Z(1+1/d)]^2,
\]
the last numerator difference being $QZ^2-d>0$. Consequently
\begin{equation}
 |y|<v\sqrt D+b+v+1
 <cT((a+2)I+1)+c+RD+1.
 \label{eq:eleven-ybound}
\end{equation}
Set $\rho=(1+1/Y)/E$. Then $c/\ell<\rho$ and $(a+3)\rho<2$.
The exact identity
\[
 (a+2)^2\rho=E+2Y+X+6+\frac YX+
 \frac5X+\frac4Y+\frac8E+\frac4{EY}
\]
has its last four terms less than $1$. As $Y/X=s_0q^2/w\le s_0q^2$
and $X\ge16$, it is less than $C_*-2$. Divide
\eqref{eq:eleven-ybound} by $\sqrt A>\ell$, and use $(c+1)/\ell<2$:
\[
 \frac{|y|}{\sqrt A}
 <2TI+R(C_*-2)i^2c^3+2\le M.
\]
No lower bound on $k$ beyond $k\ge2$ was needed.
\end{proof}

The same outer bounds give $G>A$ and $b_0>\sqrt A$: use
$\Delta c^2>k^2Y^2(E+Y)^2>A-1$ and
$b_0=\ell+E\eta+Yc+X>\ell+k>\sqrt A$.
Every root of the inner polynomial has the form $x=M\tau+y$, hence
$|x|<2M\sqrt A$.
If a vanishing factor in \eqref{eq:eleven-norm} had $g=-\sqrt G$, then
\[
 x=\frac{W+L\sqrt G}{S}
   =t+\frac{k-R-1}{E}
     +(M+1/H)(b_0+\sqrt G)>2M\sqrt A,
\]
because the first two terms have positive sum by $E>R+1$.
This is impossible. If $G$ were nonsquare, actual quadratic conjugation
would force this forbidden branch along with the positive branch. Thus
$g$ is a positive integer and $x=(W-Lg)/S$ is rational.

In a normal splitting field, an automorphism negating an irrational $\tau$
would give $2M\sqrt A\le|y|+|y'|<2M\sqrt A$. Thus $\tau$ is rational,
hence integral. Then $y=x-M\tau$ is rational, and
$f=(e-dy^2)/B_a$ is rational, hence integral. Finally the inherited rational
Pell saturation lemma gives $y\in\Z$: if $m\ge2$, $V\in\Z$, and
$m^2V^2-(m^2-1)y^2=1$ with $y\in\mathbb Q$, then $y\in\Z$.
It follows by clearing a reduced denominator, classifying the two resulting
integer Pell solutions, and using abscissa divisibility to force ordinate
divisibility. Here apply it with $m=\Delta I$ and $V=vf-b$.
The complete argument, including zero cases, is reproduced in the accompanying
proof file. In particular $x$ is integral. The carrier equality implies
\begin{equation}
 E\mid H(k-R-1).
 \label{eq:eleven-weakdiv}
\end{equation}
At this stage it would be invalid to cancel $H$.

\subsection{Recover the main index before the index quotient}

Let $A_0=a+2$, and write $\chi_{A_0},\psi_{A_0}$ for the usual Pell
sequences. The fundamental positive integer-coordinate solution for
$A_0^2-1$ is $(A_0,1)$. Main and normalized strong norm recovery gives
\[
 g=\chi_{A_0}(p),\quad c=\psi_{A_0}(p),\quad
 |f|=\chi_{A_0}(m),\quad I=\psi_{A_0}(m),\qquad p,m\ge1.
\]
The standard strong-divisibility and composition identities now give
\begin{equation}
 p\mid m,\qquad c\mid m,\qquad m\ge pc>2p.
 \label{eq:eleven-rank}
\end{equation}
Here is the direct argument, to make its independence from the deleted
equation explicit. Since $c^2\mid I$, strong divisibility first gives
$p\mid m$. Write $m=pj$. Then
$\psi_{A_0}(pj)=c\psi_{\chi_{A_0}(p)}(j)$, while modulo $c$ the
second factor is $j\chi_{A_0}(p)^{j-1}$. Since
$\chi_{A_0}(p)^2\equiv1\pmod c$, we get $c\mid j$. No ordinary-strong
large-$c$ rank lemma is necessary. Also $c>2Y>6R$, and $p\ge2$ with
$c\ge(2A_0-1)^{p-1}>2p$.

The inherited signed auxiliary-index argument can therefore be applied
before the main congruence modulo $H$. For clarity, its essential steps are
included. Put $S_0=\Delta I$ and $V=c(Tf-1)-Rf^2$. The auxiliary Pell
equation shows that
\[
 S_0|V|=\chi_{S_0}(j_0),\qquad |y|=\psi_{S_0}(j_0)
\]
for a positive odd $j_0=2r+1$. The usual odd-index polynomial $Q_r$ has
\[
 \chi_{S_0}(2r+1)=S_0Q_r(S_0^2),\quad
 Q_r(1-A_0^2)=(-1)^r\psi_{A_0}(2r+1),\quad
 Q_r(0)=(-1)^r(2r+1).
\]
Because $Q=\Delta(f^2-1)\equiv1-A_0^2\pmod{|f|}$ and
$V\equiv-c\pmod{|f|}$, these give
$\psi_{A_0}(j_0)\equiv\pm c\pmod{|f|}$. Squaring and doubling yields
\[
 \chi_{A_0}(2j_0)\equiv\chi_{A_0}(2p)\pmod{\chi_{A_0}(m)}.
\]
Signed index reduction, using $0<2p<m$, gives $j_0\equiv\pm p\pmod m$,
hence modulo $c$. This follows by reducing modulo $2m$, reflecting indices
to $[0,m]$, and using strict growth and
$2\chi_{A_0}(m-1)<\chi_{A_0}(m)$ to exclude a negative representative.
Meanwhile $Q\equiv0\pmod{c^4}$ and $V\equiv-R\pmod c$ give
$j_0\equiv\pm R\pmod c$. Since $p,R\in(0,c/2)$, one obtains
\begin{equation}
 p=R.
 \label{eq:eleven-pR}
\end{equation}

\subsection{A short index range restores the full divisibility}

The first norm has fundamental unit $(2EY+1,2)$, so for $P_0=2EY+1$,
\[
 |\tau|=\chi_{P_0}(n),\qquad k=2\psi_{P_0}(n),\qquad n\ge1.
\]
Because $P_0>A_0$ and $c>k$, we have $n<p=R$. Put
$g_0=\gcd(E,H)$ and $\mu=E/g_0$. Since $H=4XY+4Y+3$,
\begin{equation}
 \gcd(Y,H)=\gcd(Y,3)\le3,\qquad
 g_0\mid X\gcd(Y,H),\qquad \mu\ge Y/3>R.
 \label{eq:eleven-modulus}
\end{equation}
Equation \eqref{eq:eleven-weakdiv} gives $k\equiv R+1\pmod\mu$.
As $P_0\equiv1\pmod E$, the Pell recurrence also gives
$k\equiv2n\pmod E$. Hence $2n\equiv R+1\pmod\mu$. But
$1\le n\le R-1$ gives $|2n-R-1|<R<\mu$. Therefore
\[
 2n=R+1,\qquad E\mid(k-R-1).
\]
Set $h=(k-R-1)/E$. It is integral and positive: $R>3q\ge48$ gives
$n\ge25$ and $k=2\psi_{P_0}(n)>2n=R+1$.

Divide the carrier identity by $E$ to get
\[
 L_0b_0+H(t+h)=L_0g+H(M\tau+y).
\]
Since $L_0\equiv1\pmod H$, we obtain $g\equiv b_0\pmod H$.
Thus $\gamma=(g-b_0)/H$ is an integer. Moreover $n\ge2$ gives
$k\ge4P_0>X$, hence $c>X$, and
$g-b_0>(a+1)c-ac-X=c-X>0$. So $\gamma>0$.
All original normalized compiler equations are now restored, including
$k-hE-R=1$ and the main congruence. The signed-quotient soundness theorem
handles a possible negative $f$ by changing $f,T$ together, and extracts a
blank K215 halt. The order of recovery above avoids every use of $h$ before
its reconstruction.

\subsection{Completeness and exact degree}

For a canonical positive compiler witness the inherited theory gives
$X=2^R$, $R>3q$, and $Y>2^R$. Since $q\ge16$ and
$2^{3q}>3q^4$, its old coordinate $s_{\rm old}=Y/q^3$ satisfies
$s_{\rm old}>3q$. Hence $s=s_{\rm old}-3q>0$ is legitimate. Define
\[
 t=L_0\gamma+M\tau+y-h.
\]
This is positive because $h<k/E$ and $\tau>\ell=EkY>k/E$.
It makes a coherent factor in \eqref{eq:eleven-norm} vanish and preserves
$N_t=1$. This proves both directions of existence equivalence.
The shifted system uses
\[
 (K_0+w)C_w=F+z(q-1),
\]
with $z_{\rm old}=z+1$. Canonical witnesses have $w=2^R/q>q>F$ and
$C_w>0$, so $z_{\rm old}\ge2$. This saves one gate while preserving
positivity and existence.

For the degree calculation, let $\mathcal N_{\rm old}$ be the predecessor's
sharp half-carrier norm, of degree $1056$. In the localization at $E$,
\begin{equation}
 \mathcal N=E^{16}\mathcal N_{\rm old}
       \left(t+\frac{k-R-1}{E}\right),\qquad s_{\rm old}=s+3q.
 \label{eq:eleven-degreeidentity}
\end{equation}
Each of the eight factors gains $E^2$, and its degree in the old carrier
is at most $16$, so the denominators clear. The rational shift has total
degree $-2$, below $\deg t=1$, and cannot change a highest homogeneous part.
Writing $q_0=\beta J$, $k_0=\eta+\zeta$, and $s_{00}=s+3q_0$, the
predecessor's proved leading form gives
\[
 \mathcal N_{1152}
 =2^{46}i^{48}k_0^{96}w^{90}s_{00}^{192}
 q_0^{710}(q_0-F)^{16}\ne0.
\]
Thus $\deg\mathcal N=1056+16\cdot6=1152$. The transport leading form is
$w(q_0-F-\alpha)-zq_0\ne0$, proving exact product degree $1154$ and
system residual degrees $2,1152$, uniformly for the authentic $\beta>0$.

The delivered circuit has $75$ multiplications and $48$ additions or
subtractions for the product. Independent checks include $90$ formal norm
products modulo three primes, $77$ carrier-denominator identities, exact
integer and rational fixtures, $500$ parity fixtures, and two complete
univariate polynomial calculations of degree and leading coefficient.
Another $320$ bound fixtures deliberately satisfy $k\le E$, testing the
hypothesis removed from the predecessor argument. The four actual fixed
payloads are hashed and used in $40$ additional modular norm checks.
These diagnostics support the transcription; the unbounded proof above and
the inherited compiler calibration establish the theorem in ZFC.

\subsection{An optional degree-1140 companion}
\label{sec:eleven-optional}

A further positive translation of the same eleven supplied coordinates gives
a second Pareto point. Retain all definitions in Section~\ref{sec:eleven},
put
\[
 \sigma=M(\ell+2c-X)-b,
 \qquad
 W=Lb_0+H\bigl(E(t+\sigma)+k-R-1\bigr),
\]
and change nothing else. The product now has exact degree $1140$ and $129$
binary operations; the shifted two-equation version has residual degrees
$2,1138$ and $124$ operations. Thus the main $123$-operation construction
remains the less expensive one. The optional circuits are separate files.

\paragraph{Soundness of the positive translation.}
The unchanged parity argument gives $N_t=1$ at any product root and hence
all outer bounds. Write $C=C_*$ and $B=\ell+2c-X>0$. Using
$X/q=w\ge1$, direct expansion gives
\begin{align*}
 C\ell-\Delta c
 &=EY\{\eta(X-6-1/q)+\zeta(2X-4)\}-4E\eta\\
 &\hspace{1cm} +(w-1)Y^2c-(4Y+3)c\\
 &\ge EY\{\eta(X-7-1/q)+\zeta(2X-4)\}-(4Y+3)c.
\end{align*}
Since $X,q\ge16$ and $Y\ge4$, both bracket coefficients are at least
$X/2$. As $c<k(Y+1)$ and $(4Y+3)(Y+1)<8Y^2$,
\[
 C\ell-\Delta c>kY^2(X^2/2-8)\ge(k/3)X^2Y^2.
\]
Also $C=XY+2X+2Y+Y/q<3XY/2$. Consequently
\[
 CB-\Delta c>C\ell-\Delta c-CX
 >X^2Y(kY/3-3/2)>1.
\]
It follows that
\[
 \sigma=2TIB+Ri^2c^3(CB-\Delta c)-c-R>0.
\]
Every new positive solution therefore gives a positive old carrier
$t_{\rm old}=t+\sigma$, and the main soundness proof applies directly.

\paragraph{Completeness.}
Subtracting $\sigma$ from the canonical old carrier gives exactly
\[
 t=\gamma+M\{(\tau-\ell)-((a+2)c-g)\}+y+b-h.
\]
The two root differences obey
\[
 \tau-\ell=\frac{\ell k+1}{\tau+\ell}>\frac{k}{3},
 \qquad
 (a+2)c-g=\frac{c^2-1}{(a+2)c+g}
 <\frac{c}{a+2}<\frac{2k}{X+1}<\frac{k}{3}.
\]
Here $\ell>k$, $\tau<\ell+k$, and $c<k(Y+1)$ are enough.
Since $b>h$, the new $t$ is positive and completeness follows.

\paragraph{Exact degree.}
In the predecessor's sixteen-linear-factor expansion the prefactor
$d^8H^{16}$ has degree $464$, and the linear factors have degrees
four times $35$, twice $36$, and ten times $38$. The translation has degree
$35$. For exactly two factors, namely the choices
$g=\sqrt G$, $\tau=\sqrt A$, and $y=-(b-vf)+O(\lambda^{-8})$,
the translated factor becomes
\[
 t+M\{\ell-\sqrt A+(a+2)c-\sqrt G\}-\gamma-vf
    +O(\lambda^{-8}).
\]
Because $\sqrt A-\ell=k/2+O(\lambda^{-9})$, while
$(a+2)c-\sqrt G$ has degree $-1$, its leading term is $-M_0k_0/2$
of degree $28$. The opposite two degree-$35$ factors retain their degree
with twice their original leading term; the other factors are unchanged
in highest degree. The translated predecessor norm therefore has degree
\[
 464+2\cdot28+2\cdot35+2\cdot36+10\cdot38=1042.
\]
The factor $E^{16}$ from index absorption adds $96$; the residual rational
carrier correction has degree $-2$ and affects none of these leaders.
Thus the new norm has degree $1138$ and leading form
\[
 2^{46}i^{48}k_0^{96}w^{90}s_{00}^{188}
 q_0^{700}(q_0-F)^{16}.
\]
The degree-$2$ transport factor yields degree $1140$ for the single equation.
Two complete univariate polynomial checks over different finite fields agree
with this form, and $60$ independent formal conjugate-product checks verify
the optional circuit's norm transcription.

\section{Complemented whole-label clauses reduce the fixed coefficients}
\label{sec:complemented-compiler}

The improvements in the number of unknowns can use the previous coefficient
tuple.  A separate, compatible reduction makes the fixed native compiler
smaller.  We keep the ordered 943-tile language, native positions, transport
geometry, and shared-anchor synchronization of package~36 at the pinned
repository commit~\cite{repo}.  We replace its bitwise label clauses by
whole-label clauses and recompute the complete tuple and rank controller.
The resulting base is
\[
 \rho=2^{77},\qquad L=8807,\qquad d=77L=678139,\qquad B=\rho^L.
\]
The four native ports occupy $339{,}476$ raw bytes, versus $506{,}664$ in
package~36.  The largest port has $681{,}374$ bits, versus $1{,}016{,}946$.
Both reductions are approximately $33.0\%$.  These figures measure the
native ports, not the coefficients of an expanded final polynomial.

\subsection{One guard bit per entire label}

\begin{lemma}[Complemented label field]
Let $0\le x<2^{\ell-1}$ and let $v$ be an $\ell$-bit copied word.
Give a selected tile coefficient $x+1$, its copied bits their ordinary
binary weights, and the combined field $\ell+1$ bits.  Testing that the
lowest $\ell$ bits vanish forces the copied word to be
\[
 \operatorname{code}(x)=2^\ell-1-x
\]
when one tile is selected, and forces it to be zero when none is selected.
\end{lemma}

\begin{proof}
With no tile the field is $v\in[0,2^\ell-1]$.  With one tile it is
$x+1+v$, strictly between $0$ and $2^{\ell+1}$.  Its only possible multiple
of $2^\ell$ is therefore $2^\ell$.  The code is injective and its highest
bit is one.  The extra bit prevents overflow into the next field.
\end{proof}

The old clauses assigned a separate radix-four field to every copied bit.
For the four label widths $9,10,9,10$, this required $76$ payload bits.
The new fields require $10+11+10+11=42$ bits.  The three two-bit anchor
fields remain, and are placed first to reduce the carry bound.

The selector is decoded before this lemma is used.  At most $943$ tile
bits can be selected, while every other clause coefficient is divisible
by $1024$.  Thus the lowest ten bits are the exact number $n$ of selected
tiles.  Disjointness with the mask $1022$ forces $n\in\{0,1\}$.
This reasoning does not assume carry-free label fields before selection.

\subsection{A complete finite recipe}

The tile slots are west, north, east, south.  Their original labels fit
below $2^{\ell_s-1}$.  Use the following native offsets and clause offsets:
\[
\begin{array}{c|rrrr}
 s&\mathrm W&\mathrm N&\mathrm E&\mathrm S\\\hline
 o_s&1023&1013&1003&983\\
 \ell_s&9&10&9&10\\
 p_s&16&47&37&26
\end{array}
\]
The three native anchor positions are $1095,2195,3296$, with clause
offsets $10,12,14$.  Define the unchanged support
\[
 \mathcal E=\{0,\ldots,942\}
 \cup\{o_s+j:0\le j<\ell_s\}
 \cup\{1063,1064,1095,2195,3296\},
 \qquad |\mathcal E|=986.
\]
For tile $i$ with original labels $x_{i,s}$, put
\begin{align*}
 c_i&=1+\sum_s(x_{i,s}+1)2^{p_s}+2^{10}+2^{12}+2^{14},\\
 c_{o_s+j}&=2^{p_s+j},\qquad
 (c_{1095},c_{2195},c_{3296})=(2^{10},2^{12},2^{14}),\\
 c_{1063}&=c_{1064}=0,\\
 \mu&=1022+2^{10}+2^{12}+2^{14}
       +\sum_s(2^{\ell_s}-1)2^{p_s}.
\end{align*}
The clause mask has $57$ bits and population $50$.  All of the following
powers and sums define fixed integers:
\begin{align*}
 Z_0&=(2^{75}-1)\sum_{e=8792}^{8802}\rho^e
                      +(2^{39}-1)\rho^{8803},\\
 D_C&=\rho^{30}+\rho^{62}+\rho^{1100}+\rho^{2201}
      +\sum_{e\in\mathcal E}c_e
             (\rho^{4393-e}+\rho^{5495-e}),\\
 D_R&=\rho^{42},\\
 M_C&=B-1-\sum_{e\in\mathcal E\setminus\{1\}}\rho^e
                 -2\rho^{1063},\\
 M_{F,0}&=4+\mu(\rho^{4393}+\rho^{5495})
       +\sum_{j=0}^{9}\rho^{1013+j}
       +\sum_{j=0}^{8}\rho^{1065+j}
       +\rho^{2195}+\rho^{3296}+Z_0.
\end{align*}
The four ports, always used together, are
\[
 \boxed{\beta=B-1,\quad K_0=D_C+BD_R,\quad M_C,
                  \quad M_F=M_{F,0}+B-1.}
\]
In particular $M_F$ is the shifted source mask.  The unshifted mask
$M_{F,0}$ is not an interchangeable port.

\subsection{Why the native compiler still applies}

The scalar native-basis identities are unchanged: at targets $1013+j$,
$0\le j<10$, the current and temporal contributions are from $983+j$
and $1013+j$; at $1065+j$, $0\le j<9$, the current and spatial
contributions are from $1003+j$ and $1023+j$; each tested anchor
$2195,3296$ reads the current source $1095$ and its own temporal bit.
At each clause target $T\in\{4393,5495\}$, the coefficient of
$\rho^{T-e}$ in $D_C$ is precisely $c_e$.  The finite auditor checks all
$66{,}062$ identities, including the absence of every other contribution.

The occupied high bits and injective codes are retained, so the previous
matching and occupation-propagation proof applies.  The clean-band and
anchor intersections remain
\begin{align*}
 ((4393-\mathcal E)\bmod L)
       \cap((5495-\mathcal E)\bmod L)&=\varnothing,\\
 ((2195-\mathcal E)\bmod L)
       \cap((3296-\mathcal E)\bmod L)&=\{0\}.
\end{align*}
Indeed, $1102$ is absent from $\mathcal E-\mathcal E$, whereas the
difference $1101$ occurs only between the two tested anchors.  The
inequality $L>2\max\mathcal E$ excludes new circular cases.  Start first
forces the source anchor, then both tested anchors; the old binary and
whole-radix synchronization arguments recover alignment before global
occupancy is used.

For the inherited conservative pre-decoding bound, put
\[
 \mathcal M=988\left(2\sum_{e\in\mathcal E}c_e+6\right).
\]
Exact integer arithmetic verifies
\[
 \max\{4\mathcal M+8,\,2\mu+4,\,2(2^{75}-1)+4,\,16\}<2^{77}.
\]
Thus the uniform carry bounds still hold.  The largest unshifted current
exponent is $5495+3296=8791$, so the padded positions $8792$ through
$8803$ are identically zero on intended aligned cells.  The padding
contributes $11\cdot75+39=864$ tests and the original mask contributes
$122$.  Consequently
\[
 \operatorname{pc}(M_{F,0})=986,\qquad
 \operatorname{pc}(M_C)=d-986,\qquad
 M_C\equiv2\pmod4,\quad M_{F,0}\equiv4\pmod8,
\]
with $0<M_C,M_{F,0}<B-1$.  These are identities of the actual materialized
ports, not only symbolic population counts.

\subsection{An explicit controller for the actual packed rank}

Changing the radix requires a new completeness certificate.  Set
\[
 \lambda=56611,\qquad k=\lambda-1=56610,
 \qquad d=7\cdot11\cdot8807.
\]
All four displayed primes are certified by trial division.  With
$H_0=1+BD_C+B^2(D_R+1)$, exact modular arithmetic gives
\[
\begin{array}{c|rrrr}
 p&7&11&8807&56611\\\hline
 \operatorname{ord}_p(B)&3&10&629&6290\\
 (B^2-1)\bmod p&3&2&3507&20663\\
 H_0\bmod p&6&1&3377&14311
\end{array}
\]
Also $\operatorname{ord}_d(B)=18870\mid k$ and
\[
 B^k\bmod\lambda^2=1810929280\ne1,
 \qquad B^k\equiv1\pmod\lambda.
\]
Hence $v_\lambda(B^k-1)=1$.

Pad a canonical accepting rectangle to independently large powers of
$\lambda$ in both directions, as allowed by the inherited language
theorem.  Write its width and height as $w_c,h_c$, and set
$N=w_ch_c=\lambda^n$, $q=B^N$, $T=N/\lambda$, $C_0=d\lambda$.
Choose $N>d\lambda^2$ and $h_c>\lambda$.
The exact change in actual packed rank caused by the safe dummy at cell
$j$ remains
\[
 \delta R=-\Gamma B^j,\qquad
 \Gamma=\rho^{1063}(q^2-1)
          [1+q(D_C+BD_R+B^{w_c})].
\]
The unit table proves $\gcd(\Gamma,dN)=1$.

Binomial lifting and the order certificates imply that the first $T$
powers of $B^k$ are precisely the residues $1+C_0z$ modulo $C_0T$.
Every residue modulo $dN$ is a Boolean subset sum of
\[
 \{B^{2+kj}:0\le j<T\}\cup\{B^3\}.
\]
To see this, normalize a target by $B^{-2}$ and subtract either zero or
$B$ to make it nonzero modulo $\lambda$.  Its least residue $a$ modulo
$C_0$ satisfies $1\le a<C_0<T$ and $\gcd(a,T)=1$.  Write the target as
$a+C_0u$ and solve
\[
 az_0+\frac{a(a-1)}2\equiv u\pmod T.
\]
The $a$ distinct orbit entries with additive coordinates
$z_0,z_0+1,\ldots,z_0+a-1$ have the desired sum.  Include $B^3$ exactly
when the initial subtraction was used.

The largest selectable cell is at most $N-N/\lambda-\lambda+3$.
Since $w_c<N/\lambda$, its required spatial and one-cell translates
remain below $N$ and avoid the origin.  Thus the inherited no-wrap
argument applies.  Choosing the target
$(R_0-dw_c)\Gamma^{-1}$ tunes the actual rank to
$R\equiv dw_c\pmod{dN}$, as required by the helical compiler.

The unchanged mask residues give canonical $R\equiv3\pmod4$.
Since $d,N$ are odd, the canonical dyadic arithmetic scale is square;
the standard choice of strong Pell index preserves the odd normalized
multiplier, and the population identity preserves the odd scale in
$Y=sq^3$.  Thus signed adapters using these simultaneous promises remain
compatible.  Soundness does not restrict arbitrary solutions to
prime-power rectangles; this choice is only used in completeness.

The scripts \texttt{build\_complemented.py} and
\texttt{audit\_complemented.py} reproduce the integer ports and the finite
certificates.  The second script does not import the first and reconstructs
all four full integers independently.  The complete universal adapter
argument and dependency boundary are retained in
\texttt{compiler/PROOF.md}.  Rebinding a source polynomial changes no
variable count, degree, or gate count; any derived compression constants
must nevertheless be recomputed from the newly bound source.

\section{Reproduction, evidence, and remaining work}
\label{sec:verification}

\subsection{What is fully specified}

The four fixed native integers are materialized in
\texttt{compiler/coefficients/}; the binding manifest records their
unsigned little-endian raw hashes, compressed hashes, and bit lengths.
The native generator starts from a literal ordered 943-tile table and
the literal K215 transitions. It does not execute an unbounded machine
or search for a refutation. The new nine-variable $\kappa$ and
ten-integer $\beta_{\mathrm S}$ are also materialized, and every
other fixed leaf is an explicit finite integer polynomial in those
ports. Historical binding metadata in inherited source copies is
expressly inactive.

The large final polynomials are supplied as division-free directed
acyclic arithmetic circuits. This representation uniquely determines
each coefficient of a conventional integer polynomial, and permits
evaluation without printing all expanded monomials. In particular the
864-operation figure does not charge repeated shared expressions
multiple times, and it does charge all nontrivial scalar multiplications
and fixed powers after expansion. This convention is the same one
used for the repository comparisons.

\begin{longtable}{@{}p{0.27\linewidth}p{0.67\linewidth}@{}}
\toprule
Path & Purpose\\
\midrule
\endhead
\texttt{report.tex}, \texttt{sections/} & Complete manuscript source.\\
\texttt{report.pdf} & Compiled and visually inspected report.\\
\texttt{compiler/} & New native ports; builders, independent auditor,
literal data, unbounded adapter proof, and inherited notices.\\
\texttt{nine/} & Source kernel, fixed recipes, natural and positive
circuits, cancellation certificate, and independent bootstrap review.\\
\texttt{signed10/} & Ten-integer circuit, source proof, Sun interface,
finite checks, and pinned inherited finalizer.\\
\texttt{eleven/} & Main and optional eleven-positive circuits, complete
proofs, root bounds, and independent mathematical review.\\
\texttt{provenance/} & Pinned endpoint source and machine identity,
with a manifest explaining the inherited logical boundary.\\
\texttt{reproduce.py} & Runs all finite builders and auditors using
only the Python standard library.\\
\texttt{Makefile} & Convenience targets for verification and PDF build.\\
\texttt{MANIFEST.sha256} & Hashes of archive contents, excluding itself.\\
\bottomrule
\end{longtable}

From an extracted archive, run
\begin{verbatim}
python3 reproduce.py
make pdf
\end{verbatim}
Python 3.10 or later is sufficient for the finite arithmetic checks;
no third-party Python packages are required. The PDF build uses
\texttt{pdflatex} and standard TeX packages. The finite computations
need no network connection or external repository checkout. The
included Lean source snapshots identify imported theorems; they are
not a self-contained replacement for the entire upstream Lean
development.

\subsection{The executed finite certificates}

The compiler has a separate auditor that does not import its builder.
It reconstructs all four full integers and checks the label ranges,
all 66,062 scalar native-basis identities, population balance,
padding safety, and the residue orders used by the rank controller.
The builder additionally checks 1,157,632 whole-label decoder cases.
The unbounded
controller argument is proved in the manuscript and companion proof.

The nine-variable verifier independently reconstructs all 136 weighted
coefficient records using values, gradients, and quadratic parts of
the parent residuals at the all-one vector. It checks all 528 pair
sums and all strict coefficient bounds, validates actual native
payloads and the additional source constant, and compares the literal
formula, both domain versions, and the cancellation circuit in 450
finite-field evaluations using the actual fixed tuple. Its exact
symbolic diagonal audit yields the nonzero leading term
$2^{196}\kappa^{83088652}$. It also checks 768 signed-majorant and
telescoping fixtures and an exact 72-case small packing analogue with
invalid-digit controls. These latter examples detect transcription
errors; the universal source proof does not rest on their coverage.

The ten-integer checker independently derives the quadratic source,
verifies the full fixed tuple and the even source exponent, expands
the three-root norm identity formally, and compares 1620 named
registers over three finite fields. Its actual-instance diagonal
leading coefficients are nonzero modulo 101, 1009, and 1,000,003;
the leading-degree structure therefore survives the fixed coefficient
specialization. The proof-level specialization of Sun's interface is
stated separately from these algebraic tests.

The main eleven-positive checker compares 90 independently multiplied
formal conjugate products, 77 carrier-denominator identities, exact
integer and rational fixtures, 500 parity fixtures, and two complete
univariate coefficient arrays for the degree calculation. Another
320 bound fixtures deliberately have $k\le E$, so they exercise the
case in which the predecessor's deleted index condition cannot be
used. Forty additional norm checks use the actual coefficient tuple.
The optional translated construction has two further complete degree
checks and 60 independent norm products. The ordinary proofs and
independent review notes cover both positive-domain directions.

\subsection{The dependency boundary}

There are four logically different layers.
First, the finite source data and scalar identities are completely
reproducible inside this archive. Second, the new reductions---quadratic
batching, its signed bootstrap, index absorption, positive translation,
and complemented label clauses---are ordinary mathematical proofs
given here and in the accompanying proof files. Third, the normalized
native/Pell compiler, canonical completeness choices, and literal K215
logical calibration are inherited from the pinned repository.
Fourth, the unrestricted-integer companion uses the explicitly named
theorems of the external Sun papers. These sources have been identified
and checked for the contracts used, but are not re-proved in a new
proof-assistant kernel in this work.

The ordinary proofs use finite integer arithmetic, induction,
elementary algebra, and the stated inherited arithmetic and compiler
theorems. Thus the displayed equivalences are formalizable in ZFC.
That assertion is different from claiming that an automated checker
has already accepted the complete new theorem. The source files mark
this distinction consistently.

\subsection{Promising next steps}

The structural gain is already achieved without optimizing the Sidon
set. A smaller 32-position set, or a source-aware set that separates
only the needed coefficient matches, would reduce the two large-degree
companions further. Such a replacement needs its own exact collision
certificate and must retain the complement positions needed for all
target tests. Merely assigning the 136 existing monomials distinct
positions is insufficient: products with the other monomials in the
digit-code square must also be excluded.

The most direct formalization target is the general simultaneous
quadratic extraction and carry decoder. It is a short independent
lemma between the existing quadratic source and the already
parameterized ratio theorem. The signed bootstrap should be included
in that formalization, since a positive-source assumption on arbitrary
outer inputs would be false. For the eleven-variable construction,
the central new formal lemma is the reordered Pell-index recovery
followed by the short-modulus argument; the proof should not import
$k>E$ from the old twelve-variable formulation.

The present results do not settle an eight-natural-variable or
nine-unrestricted-integer construction for this consistency predicate.
They do provide a large and reproducible improvement at the existing
few-variable bounds, a new moderate-degree eleven-positive form, and
a smaller coherent coefficient compiler. Those are separate,
verifiable advances over the pinned report family.

\clearpage
\begin{thebibliography}{9}
\raggedright
\bibitem{repo}
V. Reshetnikov and contributors, \emph{ProveIt}, repository revision
\repoPin, accessed 23 October 2026.
Report family \texttt{con-zfc-diophantine-forms}, especially packages
30, 31, 34--37; source and inherited notices are recorded in the archive.
\href{https://github.com/VladimirReshetnikov/ProveIt/tree/22e75dd6b87b1b199bb08c778d7cf075ae5aae08/SetTheory/Cardinals/docs/reports/hilbert-tenth-problem/con-zfc-diophantine-forms}{Pinned report directory and source archive}.
\bibitem{jones1982}
J. P. Jones, \emph{Universal Diophantine equation}, Journal of Symbolic
Logic \textbf{47} (1982), no. 3, 549--571.
\url{https://doi.org/10.2307/2273588}.
The parameterized carry and ratio interfaces used here are present in
the repository's \texttt{Paper1982/Carries.lean} and
\texttt{Paper1982/NineRatio.lean}.
\bibitem{sun2021}
Z.-W. Sun, \emph{Further results on Hilbert's Tenth Problem}, Science
China Mathematics \textbf{64} (2021), 281--306.
\url{https://arxiv.org/abs/1704.03504}, version 7.
\bibitem{sun2026}
Z.-W. Sun, \emph{Ten unknowns for Hilbert's tenth problem over the
integers}, preprint, arXiv:2609.01594, version 2 (2026).
\url{https://arxiv.org/html/2609.01594v2}.
\end{thebibliography}
\end{document}
